\chapter{Vis\~ao geral do sistema}
\label{cap:visao}

\section{Resumo executivo}

O Rust VGA Kernel \'e um execut\'avel ELF x86\_64 carregado pelo GRUB. O GRUB
usa o protocolo Multiboot2 para transferir o controle a um pequeno \emph{stub}
em linguagem de montagem, ligado em 1 MiB. Esse stub constr\'oi tabelas de
p\'agina, ativa PAE e modo longo e salta para o c\'odigo Rust, ligado na metade
alta em \code{0xFFFFFFFF80000000}.

O c\'odigo Rust inicializa, nesta ordem, o renderizador VGA, a serial COM1, a
camada TTY com oito consoles virtuais, o alocador de frames, o heap do kernel,
o ramfs e o VFS, a GDT com ring 3/TSS, o teclado PS/2 e a IDT com PIC/PIT. Em
seguida liga as interrup\c{c}\~oes, cria a tabela de tarefas com o shell como
PID 1, imprime um \emph{banner}, valida a assinatura Multiboot2 e entra no
shell, que exibe o prompt \code{rinux>} no console ativo e processa comandos
lidos do teclado ou da serial.

O sistema inteiro cabe em cerca de 250 KiB de ELF e 6.797 linhas de Rust
(5.365 linhas efetivas de c\'odigo, sem coment\'arios e linhas em branco).
Inclui um heap \emph{first-fit} atr\'as de \code{\#[global\_allocator]}, um VFS
com dezesseis descritores sobre o ramfs e um \code{/proc} sint\'etico, um
escalonador \emph{round-robin} preemptivo por timer com \code{sleep}
bloqueante (\code{schedule\_timeout}) e filas de espera por evento
(\code{wait\_event}/\code{wake\_up}), e onze chamadas de sistema. Ainda
n\~ao h\'a \code{fork}/\code{exec}, espa\c{c}o de endere\c{c}os por processo
nem sistema de arquivos em disco.

\section{Funcionalidades}

\begin{table}[H]
\centering
\caption{Funcionalidades implementadas e seu estado.}
\label{tab:funcionalidades}
\begin{tabularx}{\linewidth}{@{}lXc@{}}
\toprule
\textbf{\'Area} & \textbf{Funcionalidade} & \textbf{Estado} \\
\midrule
Boot & Cabe\c{c}alho Multiboot2 v\'alido & Completo \\
Boot & Transi\c{c}\~ao 32 $\rightarrow$ 64 bits & Completo \\
Mem\'oria & PML4/PDPT/PD est\'aticos, 1 GiB & Completo \\
Mem\'oria & Metade alta do kernel & Completo \\
V\'ideo & Console VGA texto 80$\times$25 & Completo \\
V\'ideo & Rolagem, cursor, tab, backspace & Completo \\
V\'ideo & Oito consoles virtuais (\code{tty0}..\code{tty7}) & Completo \\
V\'ideo & Troca por Alt+F1..Alt+F8 e comando \code{tty} & Completo \\
V\'ideo & Framebuffer gr\'afico / GOP & Ausente \\
Teclado & PS/2 8042, conjunto 1, polling & Completo (subconjunto) \\
Teclado & Shift, Alt, Enter, Tab, Backspace & Completo \\
Teclado & Sele\c{c}\~ao de console com Alt+F1..Alt+F8 & Completo \\
Teclado & IRQ1, Caps Lock, AltGr, setas & Ausente \\
Serial & COM1 16550, 38400 8N1 & Completo (subconjunto) \\
Serial & Entrada pelo terminal do QEMU & Completo \\
Shell & \code{help}, \code{clear}, \code{echo} & Completo \\
Shell & \code{ls}, \code{cat}, \code{cp}, \code{mv}, \code{tree} & Completo \\
Shell & \code{info}, \code{tty}, \code{getty} & Completo \\
Shell & \code{reboot}, \code{halt} & Completo \\
Sistema & Panic handler & Completo \\
Sistema & IDT, exce\c{c}\~oes e IRQ0 (timer) & Completo (subconjunto) \\
Sistema & GDT/TSS, ring 3 e prote\c{c}\~ao U/S & Completo (demo) \\
Sistema & Syscalls \code{exit/write/read/get\_tty} & Completo (subconjunto) \\
Sistema & Syscalls de arquivo \code{open/close/read/write/lseek} & Completo (subconjunto) \\
Sistema & \code{getpid}/\code{getppid} vindos da tabela de tarefas & Completo \\
Usu\'ario & \code{getty}/login em ring 3 & Completo (demo) \\
Mem\'oria & Frames, p\'aginas 4 KiB e demand paging & Completo (subconjunto) \\
Mem\'oria & Heap do kernel, crate \code{alloc} e \code{Vec}/\code{String} & Completo (subconjunto) \\
Armazenamento & Ramfs em mem\'oria (48 n\'os, arquivos de 1 KiB) & Completo (subconjunto) \\
Armazenamento & VFS, 16 descritores e \code{lseek} & Completo (subconjunto) \\
Armazenamento & \code{/proc} sint\'etico (version, uptime, meminfo, ps) & Completo (subconjunto) \\
Processos & Tabela de tarefas, PID/PPID e estados & Completo (subconjunto) \\
Processos & Escalonador \emph{round-robin} preemptivo por timer & Completo (subconjunto) \\
Processos & \code{sleep} bloqueante (\code{schedule\_timeout}) & Completo (subconjunto) \\
Processos & Filas de espera por evento (\code{wait\_event}/\code{wake\_up}) & Completo (subconjunto) \\
Processos & \code{fork}/\code{exec}/\code{wait}, CFS/EEVDF & Ausente \\
Armazenamento & Driver de disco, mounts e permiss\~oes & Ausente \\
\bottomrule
\end{tabularx}
\end{table}

\section{Requisitos}

\subsection{Hardware}

\begin{itemize}
  \item PC compat\'ivel x86\_64 com BIOS legada ou CSM habilitado.
  \item Controlador de teclado PS/2 8042 (ou emula\c{c}\~ao equivalente).
  \item Adaptador VGA em modo texto, com buffer em \code{0xB8000}.
  \item Unidade de CD/DVD ou dispositivo equivalente para inicializar a ISO.
  \item Mem\'oria suficiente para carregar o kernel e a ISO; 128 MiB s\~ao
        suficientes no QEMU.
\end{itemize}

\subsection{Software de desenvolvimento}

\begin{itemize}
  \item Rust est\'avel com o alvo \code{x86\_64-unknown-none} instalado.
  \item \code{rustup} (recomendado; o Makefile prioriza o toolchain dele).
  \item \code{ld.lld} (LLVM linker).
  \item \code{grub-mkrescue}, \code{xorriso}, \code{libisoburn} e
        \code{mtools}.
  \item QEMU (\code{qemu-system-x86\_64}) para execu\c{c}\~ao e testes.
  \item Para esta documenta\c{c}\~ao: uma distribui\c{c}\~ao TeX Live com
        \code{pdflatex} e \code{latexmk}.
\end{itemize}

\section{Estrutura de diret\'orios}

\begin{lstlisting}[style=shell]
.
|-- Cargo.toml              metadados e perfis de compilacao
|-- Cargo.lock              lockfile (sem dependencias externas)
|-- rust-toolchain.toml     canal stable + alvo x86_64-unknown-none
|-- .cargo/config.toml      alvo padrao, linker e flags
|-- linker.ld               layout do ELF (stub baixo + kernel alto)
|-- Makefile                build do kernel, ISO e execucao no QEMU
|-- boot/grub/grub.cfg      menu Multiboot2 e console VGA
|-- src/
|   |-- main.rs             entrada Rust, macros e panic handler
|   |-- boot.rs             Multiboot2, GDT, paginacao e modo longo
|   |-- arch.rs             portas de I/O, interrupcoes, halt e reboot
|   |-- gdt.rs              GDT definitiva, TSS e pilhas de privilegio
|   |-- idt.rs              IDT, stubs, PIC, PIT e porta int 0x80
|   |-- paging.rs           frames, tabelas 4 KiB e demand paging
|   |-- heap.rs             heap do kernel e #[global_allocator]
|   |-- ramfs.rs            sistema de arquivos em memoria
|   |-- proc.rs             /proc virtual (uptime, meminfo, ps, ...)
|   |-- vfs.rs              VFS minimo e tabela de descritores
|   |-- task.rs             tarefas, escalonador, sleep e wait queues
|   |-- user.rs             ring 3, demos e getty/login
|   |-- syscall.rs          exit/write/read/open/close/lseek/getpid
|   |-- vga.rs              renderizador VGA 80x25
|   |-- tty.rs              consoles virtuais tty0..tty7
|   |-- keyboard.rs         teclado PS/2 (8042)
|   |-- serial.rs           UART 16550 / COM1
|   |-- shell.rs            prompt e comandos
|   `-- sync.rs             spinlocks (console, ramfs, heap, VFS, tarefas)
|-- docs/                   esta documentacao (LaTeX e PDF)
|-- build/                  arvore temporaria da ISO (gerada)
|-- target/                 artefatos do Cargo (gerado)
|-- vga-kernel.iso          imagem inicializavel (gerada)
`-- README.md               guia rapido de uso
\end{lstlisting}

\section{M\'etricas do c\'odigo}

\begin{table}[H]
\centering
\caption{M\'etricas dos m\'odulos Rust (contagem por linhas f\'isicas).}
\label{tab:metricas}
\begin{tabular}{@{}lrrl@{}}
\toprule
\textbf{M\'odulo} & \textbf{Linhas} & \textbf{Fun\c{c}\~oes} & \textbf{Responsabilidade} \\
\midrule
\code{boot.rs}     &  202 &  0 & Stub de boot e tabelas de p\'agina \\
\code{arch.rs}     &  101 &  7 & Portas de I/O, \code{irqsave}, halt, reboot \\
\code{gdt.rs}      &  159 &  6 & GDT, TSS e pilhas de privil\'egio \\
\code{idt.rs}      &  507 & 14 & IDT, PIC, PIT, exce\c{c}\~oes e IRQs \\
\code{paging.rs}   &  559 & 23 & Frames, tabelas e demand paging \\
\code{heap.rs}     &  270 &  9 & Heap do kernel e alocador global \\
\code{ramfs.rs}    &  534 & 31 & Sistema de arquivos em mem\'oria \\
\code{proc.rs}     &  191 &  9 & \code{/proc} sint\'etico \\
\code{vfs.rs}      &  302 & 11 & VFS e tabela de descritores \\
\code{task.rs}     &  746 & 29 & Tarefas, escalonador, sleep e wait queues \\
\code{user.rs}     &  446 &  9 & Ring 3, demos e getty \\
\code{syscall.rs}  &  244 &  9 & Syscalls de arquivo e processo \\
\code{vga.rs}      &   74 &  5 & Renderizador VGA texto \\
\code{tty.rs}      &  448 & 31 & Consoles virtuais e filas \\
\code{keyboard.rs} &  271 & 10 & Driver PS/2 8042 \\
\code{serial.rs}   &   70 &  5 & UART 16550 / COM1 \\
\code{shell.rs}    & 1468 & 61 & Linha de comando e applets \\
\code{main.rs}     &  139 &  2 & Entrada Rust, macros, panic \\
\code{sync.rs}     &   66 &  5 & Spinlocks seguros para IRQ \\
\midrule
\textbf{Total}     & \textbf{6.797} & \textbf{276} & 5.365 linhas efetivas de c\'odigo \\
\bottomrule
\end{tabular}
\end{table}

\begin{nota}[Sobre a contagem]
A contagem inclui 654 linhas de coment\'ario e 778 linhas em branco. O n\'umero
de fun\c{c}\~oes inclui m\'etodos de \code{impl}, fun\c{c}\~oes \code{const}
e fun\c{c}\~oes de trait (\code{Deref}, \code{DerefMut}, \code{Drop} e
\code{fmt::Write}); n\~ao inclui declara\c{c}\~oes \code{extern} sem corpo.
\code{boot.rs} n\~ao possui fun\c{c}\~oes Rust porque sua l\'ogica est\'a em
\code{global\_asm!}; o \emph{stub} \'e contado como bloco de montagem.
\end{nota}

\section{Fluxo resumido de execu\c{c}\~ao}

\begin{enumerate}
  \item A BIOS seleciona o CD e executa o El Torito.
  \item O GRUB carrega o ELF e l\^e o cabe\c{c}alho Multiboot2.
  \item O GRUB entra em modo protegido de 32 bits e salta para \code{\_start}.
  \item O stub salva \code{eax} e \code{ebx}, monta as p\'aginas, carrega a
        GDT e ativa o modo longo.
  \item O stub salta para \code{kernel\_main(magic, info)} na metade alta.
  \item O Rust limpa a tela VGA, configura a serial, inicializa a camada TTY
        com oito consoles virtuais, o alocador de frames, o heap do kernel, o
        ramfs e o VFS, a GDT/TSS e o teclado, e instala a IDT com PIC/PIT.
  \item O Rust liga as interrup\c{c}\~oes, cria a tabela de tarefas com o
        shell como PID 1 e o shell exibe \code{rinux>} no console ativo,
        processando comandos at\'e \code{halt}, \code{reboot} ou reset
        f\'isico.
\end{enumerate}
